\chapter{CONCLUSIONES}

El desarrollo de QuickTable demostró que es posible construir una solución tecnológica
accesible, segura y funcional para el sector gastronómico colombiano sin recurrir a
servicios en la nube ni a hardware especializado de alto costo, validando la hipótesis
central planteada en el anteproyecto. Los resultados obtenidos a lo largo de las fases
de laboratorio y de la implementación piloto en entorno real confirman el cumplimiento
de los tres objetivos específicos del proyecto.

\section*{Cumplimiento del objetivo de diseño y arquitectura}

El diseño de la arquitectura del sistema integró de forma coherente los requerimientos
funcionales, de seguridad y de escalabilidad definidos en la fase inicial. La decisión
de adoptar el patrón Razor Pages sobre .NET 8 con Entity Framework 6 y SQL Server
permitió estructurar un sistema robusto que opera íntegramente sobre intranet local,
eliminando la dependencia de conectividad externa y garantizando al establecimiento la
custodia absoluta de su información operativa y financiera. El modelo entidad-relación
diseñado separó con claridad la operación en tiempo real de la persistencia histórica,
lo que contribuyó directamente a los tiempos de respuesta inferiores a dos segundos
registrados en las pruebas de carga. La arquitectura de seguridad multicapa, compuesta
por encriptación TDE sobre los archivos físicos de la base de datos, hashing BCrypt
con factor de trabajo 12 para credenciales de usuario y cifrado AES-256-CBC para los
identificadores NFC, estableció un estándar de protección que ninguna de las soluciones
POS comerciales analizadas en el mercado colombiano contempla de forma integrada.

\section*{Cumplimiento del objetivo de implementación}

El sistema fue implementado en su totalidad, abarcando los cinco módulos operativos
(Mesero, Cocina, Caja, Administrador y TI), el módulo de seguridad por hardware 2FA
y la API REST de comunicación con el cliente Python ejecutado en la Raspberry Pi. La
interfaz de usuario, construida sobre AdminLTE con personalización CSS profunda,
demostró adaptabilidad a dispositivos heterogéneos, desde smartphones Android para
la toma de pedidos hasta televisores Smart TV para el módulo de cocina, sin dependencia
de sistemas operativos específicos. El módulo de autenticación de doble factor por
hardware, basado en la Raspberry Pi 3 con lector RFID-RC522 operando a 13,56~MHz
conforme al estándar ISO/IEC 14443, opera de forma completamente offline dentro de la
intranet, resolviendo la incompatibilidad estructural de los sistemas 2FA en la nube
con el modelo de despliegue local del sistema. La implementación del modelo de
concurrencia threading en el cliente Python eliminó el bloqueo de la interfaz gráfica
Tkinter durante las lecturas NFC, garantizando retroalimentación visual continua al
administrador durante el proceso de autenticación.

\section*{Cumplimiento del objetivo de evaluación}

Las pruebas de validación ejecutadas confirmaron la viabilidad técnica del sistema
en todas las dimensiones evaluadas. Las siete pruebas unitarias MSTest sobre los
servicios PasswordService y CryptoService registraron estado \textit{Passed} en dos
ejecuciones independientes, verificando la correcta implementación de los algoritmos
criptográficos y su comportamiento ante entradas inválidas o vacías. Las pruebas de
API con Postman validaron la coherencia semántica de los cuatro endpoints evaluados,
confirmando la operatividad del sistema de copias de seguridad con TDE integrado y el
correcto registro de asistencia del personal mediante tarjeta NFC. La evaluación de
rendimiento con Lighthouse evidenció tiempos de respuesta operativa inferiores a dos
segundos en las condiciones de intranet aislada, con puntajes de 99 a 100 en las vistas
de Login, Cocina, Caja y Administrador, y un puntaje de 87 en la vista de Mesero
atribuible al contexto de emulación de red móvil con limitaciones de procesamiento
simuladas. El análisis de vulnerabilidades con OWASP ZAP 2.17.0 sobre 15 endpoints
identificó únicamente alertas de nivel medio e informativo, concentradas en
configuraciones de política de seguridad de contenido requeridas por las librerías de
interfaz de terceros, sin detectar vulnerabilidades que comprometieran la integridad
de la base de datos ni la persistencia del sistema.

La prueba de concurrencia automatizada con Playwright, ejecutada durante siete horas
continuas con seis agentes paralelos que simularon los roles de Mesero, Cocina y Caja,
registró tasas de éxito transaccional del 100\,\% en los roles de Caja y Mesero, y del
99,75\,\% en el rol de Cocina. El único evento categorizado como error en el módulo de
cocina correspondió a una condición de carrera momentánea durante una modificación
concurrente, resuelta automáticamente por la arquitectura de Razor Pages sin propagación
al resto de los procesos del servicio. El sistema operó durante las siete horas de
ejecución sin desbordamientos de memoria ni bloqueos de recursos físicos, confirmando
la disponibilidad sostenida del servicio bajo los parámetros de infraestructura propios
de una PYME gastronómica.

La documentación técnica y de usuario generada como parte del cuarto objetivo específico,
que comprende el manual técnico de instalación y configuración del sistema y los manuales
de usuario diferenciados por rol operativo, se encuentra disponible en el Anexo~A de este
documento.

\section*{Validación en entorno real}

La implementación piloto en Gastro Uno 58 Bar durante un mes de operación continua
constituyó la validación más exigente del sistema, al exponer el software a condiciones
reales de uso por parte de personal sin formación técnica previa. Los resultados del
cuestionario de satisfacción aplicado al finalizar el período mostraron acuerdo unánime
del 100\,\% en las ocho dimensiones evaluadas: eficiencia operativa, estabilidad técnica,
usabilidad, reducción de errores, comunicación interna, control de asistencia,
autenticación de doble factor y gestión del negocio. La calificación promedio de
satisfacción general fue de 4,80 sobre 5,00, con cuatro de los cinco participantes
otorgando la puntuación máxima. Ningún participante calificó el sistema por debajo de
cuatro estrellas. La retroalimentación cualitativa recopilada, que incluyó solicitudes
de mejora sobre el módulo de asistencia, la programación de menús por día y el formato
de impresión de facturas, no señaló fallos funcionales sino oportunidades de mejora
incremental, todas documentadas en el capítulo de Trabajo Futuro.

\section*{Viabilidad económica y posicionamiento competitivo}

El análisis de viabilidad económica confirmó que QuickTable representa la alternativa
de menor costo total de propiedad disponible en el mercado colombiano para el segmento
PYME gastronómico. Frente a los sistemas POS comerciales con suscripción mensual
identificados en el mercado, cuyo costo real para el restaurante oscila entre \$105.800
y \$251.890~COP mensuales al incluir el internet requerido para su operación, QuickTable
alcanza el punto de equilibrio económico antes del décimo mes en el escenario más
conservador (frente a GastroPOS Restobar) y antes del quinto mes frente a los
competidores del segmento de gestión completa. A dos años de operación, el ahorro
acumulado del restaurante frente a Loggro Básico supera los \$3.300.000~COP. La única
dimensión en la que QuickTable no compite en igualdad de condiciones es la facturación
electrónica DIAN, ausente en la versión actual del sistema, lo que restringe su adopción
en establecimientos con obligación de emitir facturas electrónicas bajo régimen común.
Esta limitación ha sido identificada como la línea de desarrollo de mayor impacto para
versiones posteriores del sistema.

\section*{Aporte interdisciplinario}

Más allá del resultado funcional, este proyecto demostró la aplicabilidad de los
principios de la Ingeniería Electrónica, incluyendo comunicaciones inalámbricas de corto
alcance, sistemas embebidos y arquitecturas de seguridad por hardware, a la solución de
problemas reales del sector productivo nacional. La integración del módulo NFC conforme
al estándar ISO/IEC 14443 con el backend .NET 8, la gestión de concurrencia de hilos en
el cliente Python y el diseño del protocolo de intercambio de tokens de sesión de corta
duración constituyen contribuciones técnicas que trascienden el ámbito del desarrollo
puramente informático y posicionan al proyecto como un caso de convergencia efectiva
entre hardware embebido y software de nivel empresarial orientado al segmento PYME
colombiano.